\chapter{Atoms}

This chapter develops the theory of atoms in a vector lattice. An
\emph{atom} (or \emph{discrete element}) is a strictly positive vector
$a$ whose order interval $[0, a]$ collapses to the half-line
$\mathbb{R}_+ a$: every $x$ with $0 \le x \le a$ is a non-negative
scalar multiple of $a$. We record several characterisations of atoms
— the absence of two non-zero disjoint elements in $[0, a]$ and the
coincidence of the principal ideal of $a$ with the line
$\mathbb{R} a$ — together with the dichotomy that any two atoms are
either disjoint or proportional, and the fact that the band generated
by an atom is itself a projection band that coincides with
$\mathbb{R} a$. Every Archimedean vector lattice $X$ then admits a
canonical decomposition into its \emph{atomic part} — the band
generated by all atoms — and its \emph{continuous part} — the disjoint
complement of the atomic part; the two bands meet only at zero. A
vector lattice is \emph{atomic} (or \emph{discrete}) when it equals
its atomic part. The chapter culminates in Yudin's theorem: an
Archimedean vector lattice is finite-dimensional if and only if it is
vector-lattice isomorphic to $\mathbb{R}^n$ for some $n$. As a final
consequence, in an atomic Archimedean vector lattice every
non-negative element admits a coordinate expansion as the supremum of
its scalar multiples along any maximal disjoint family of atoms.

\section{Atoms}

This section introduces the notion of atom of a vector lattice and
develops its first properties. We give two characterisations of atoms
in an Archimedean vector lattice — via the absence of two disjoint
non-zero elements in the order interval $[0, a]$ and via the
coincidence of the principal ideal of $a$ with the line $\mathbb{R} a$
— show that the family of atoms is closed under positive scalar
multiplication, and prove the dichotomy that any two atoms are either
disjoint or scalar multiples of each other.

\begin{definition}
\label{def:isVLAtom}
\lean{IsVLAtom}
\leanok
\uses{def:vector-lattice}
Let $X$ be a vector lattice. An element $a \in X$ is an \emph{atom}
(or \emph{discrete element}) of $X$ when $0 < a$ and every $x \in X$
with $0 \le x \le a$ is a real scalar multiple of $a$.
\end{definition}

\begin{lemma}
\label{lem:eq-zero-or-eq-zero-of-isVLDisjoint-of-isVLAtom}
\lean{eq_zero_or_eq_zero_of_isVLDisjoint_of_isVLAtom}
\leanok
\uses{def:isVLAtom, def:isVLDisjoint, lem:inf-eq-zero-of-isVLDisjoint,
  lem:inf-smul-nonneg}
Let $X$ be a vector lattice, let $a \in X$ be an atom, and let
$x, y \in X$ with $0 \le x \le a$, $0 \le y \le a$ and $x$, $y$
disjoint. Then $x = 0$ or $y = 0$.
\end{lemma}

\begin{theorem}
\label{thm:isVLAtom-iff-no-disjoint-in-interval}
\lean{isVLAtom_iff_no_disjoint_in_interval}
\leanok
\uses{def:isVLAtom, def:isVLDisjoint, def:is-vl-archimedean,
  lem:eq-zero-or-eq-zero-of-isVLDisjoint-of-isVLAtom,
  lem:isVLDisjoint-posPart-negPart}
Let $X$ be an Archimedean vector lattice and let $a \in X$ with
$0 < a$. Then $a$ is an atom of $X$ if and only if for all
$x, y \in X$ with $0 \le x \le a$, $0 \le y \le a$ and $x$, $y$
disjoint, $x = 0$ or $y = 0$.
\end{theorem}

\begin{lemma}
\label{lem:isVLAtom-smul}
\lean{IsVLAtom.smul}
\leanok
\uses{def:isVLAtom}
Let $X$ be a vector lattice, let $a \in X$ be an atom, and let
$\lambda \in \mathbb{R}$ with $0 < \lambda$. Then $\lambda a$ is an
atom of $X$.
\end{lemma}

\begin{theorem}
\label{thm:isVLAtom-iff-principal-eq-span}
\lean{isVLAtom_iff_principal_eq_span}
\leanok
\uses{def:isVLAtom, def:is-vl-archimedean, def:orderIdeal-principal}
Let $X$ be an Archimedean vector lattice and let $a \in X$. Then
$a$ is an atom of $X$ if and only if $0 < a$ and every element of
the principal ideal $I_a$ is a real scalar multiple of $a$.
\end{theorem}

\begin{theorem}
\label{thm:isVLDisjoint-or-smul-of-isVLAtom}
\lean{isVLDisjoint_or_smul_of_isVLAtom}
\leanok
\uses{def:isVLAtom, def:isVLDisjoint, def:is-vl-archimedean,
  lem:isVLDisjoint-of-inf-eq-zero}
Let $X$ be an Archimedean vector lattice and let $a, b \in X$ be
atoms. Then either $a$ and $b$ are disjoint, or there exists
$c \in \mathbb{R}$ with $b = c\,a$.
\end{theorem}

\section{The principal band of an atom}

For an atom $a$ of an Archimedean vector lattice the band generated
by $\{a\}$ admits a particularly transparent description: it
coincides with the line $\mathbb{R} a$, and is moreover a projection
band of $X$.

\begin{theorem}
\label{thm:principalBand-eq-smul-of-isVLAtom}
\lean{principalBand_eq_smul_of_isVLAtom}
\leanok
\uses{def:isVLAtom, def:is-vl-archimedean, def:band-generated,
  thm:band-mem-generated-singleton-iff-isLUB-inf-nsmul}
Let $X$ be an Archimedean vector lattice and let $a \in X$ be an
atom. Then the band generated by $\{a\}$ equals
$\mathbb{R} a := \{c\,a : c \in \mathbb{R}\}$.
\end{theorem}

\begin{theorem}
\label{thm:exists-projectionBand-principalBand-of-isVLAtom}
\lean{exists_projectionBand_principalBand_of_isVLAtom}
\leanok
\uses{def:isVLAtom, def:is-vl-archimedean, def:band-generated,
  def:projectionBand, def:disjointComplement,
  thm:projectionBand-iff-add-disjointComplement,
  thm:principalBand-eq-smul-of-isVLAtom}
Let $X$ be an Archimedean vector lattice and let $a \in X$ be an
atom. Then the band generated by $\{a\}$ is a projection band of
$X$.
\end{theorem}

\section{Atomic and continuous parts}

Every vector lattice $X$ admits a canonical pair of bands: the
\emph{atomic part}, the band generated by all atoms of $X$, and the
\emph{continuous part}, the band of vectors disjoint from every
atom. In the Archimedean setting, the continuous part is precisely
the disjoint complement of the atomic part, and the two bands meet
only at zero.

\begin{definition}
\label{def:atomicPart}
\lean{atomicPart}
\leanok
\uses{def:isVLAtom, def:band-generated}
Let $X$ be a vector lattice. The \emph{atomic part} of $X$ is the
band generated by the set of all atoms of $X$.
\end{definition}

\begin{definition}
\label{def:continuousPart}
\lean{continuousPart}
\leanok
\uses{def:isVLAtom, def:disjointComplement, lem:band-disjointComplement}
Let $X$ be a vector lattice. The \emph{continuous part} of $X$ is
the disjoint complement of the set of all atoms of $X$ —
equivalently, the band of vectors of $X$ disjoint from every atom.
\end{definition}

\begin{lemma}
\label{lem:continuousPart-eq-disjointComplement-atomicPart}
\lean{continuousPart_eq_disjointComplement_atomicPart}
\leanok
\uses{def:atomicPart, def:continuousPart, def:is-vl-archimedean,
  thm:band-disjointComplement-disjointComplement-eq-generated,
  lem:disjointComplement-disjointComplement-disjointComplement}
Let $X$ be an Archimedean vector lattice. Then the continuous part
of $X$ equals the disjoint complement of the atomic part of $X$.
\end{lemma}

\begin{theorem}
\label{thm:atomicPart-inter-continuousPart-eq-zero}
\lean{atomicPart_inter_continuousPart_eq_zero}
\leanok
\uses{def:atomicPart, def:continuousPart, def:is-vl-archimedean,
  lem:continuousPart-eq-disjointComplement-atomicPart,
  lem:disjointComplement-inter-eq-zero}
Let $X$ be an Archimedean vector lattice and let $x \in X$. If $x$
belongs to both the atomic and the continuous parts of $X$, then
$x = 0$.
\end{theorem}

\section{Atomic vector lattices}

A vector lattice is \emph{atomic} (or \emph{discrete}) when it is
generated, as a band, by its atoms. This section introduces this
notion as a typeclass and gives three equivalent characterisations
in the Archimedean setting: triviality of the continuous part, the
existence of an atom below every strictly positive element, and the
equality of the continuous part with $\{0\}$.

\begin{definition}
\label{def:isAtomicVectorLattice}
\lean{IsAtomicVectorLattice}
\leanok
\uses{def:atomicPart}
Let $X$ be a vector lattice. We call $X$ \emph{atomic} (or
\emph{discrete}) when every $x \in X$ belongs to the atomic part of
$X$.
\end{definition}

\begin{theorem}
\label{thm:isAtomicVectorLattice-iff-continuousPart-eq-zero}
\lean{isAtomicVectorLattice_iff_continuousPart_eq_zero}
\leanok
\uses{def:isAtomicVectorLattice, def:continuousPart,
  def:is-vl-archimedean,
  thm:atomicPart-inter-continuousPart-eq-zero,
  thm:band-eq-disjointComplement-disjointComplement,
  lem:continuousPart-eq-disjointComplement-atomicPart}
Let $X$ be an Archimedean vector lattice. Then $X$ is atomic if and
only if the only $x \in X$ in the continuous part of $X$ is $x = 0$.
\end{theorem}

\begin{theorem}
\label{thm:isAtomicVectorLattice-iff-forall-pos-dominates-atom}
\lean{isAtomicVectorLattice_iff_forall_pos_dominates_atom}
\leanok
\uses{def:isAtomicVectorLattice, def:isVLAtom,
  def:is-vl-archimedean,
  thm:isAtomicVectorLattice-iff-continuousPart-eq-zero,
  lem:abs-eq-zero-iff-zero}
Let $X$ be an Archimedean vector lattice. Then $X$ is atomic if and
only if every $x \in X$ with $0 < x$ dominates an atom: there
exists an atom $a \in X$ with $a \le x$.
\end{theorem}

\begin{theorem}
\label{thm:isAtomicVectorLattice-iff-continuousPart-trivial}
\lean{isAtomicVectorLattice_iff_continuousPart_trivial}
\leanok
\uses{def:isAtomicVectorLattice, def:continuousPart,
  def:is-vl-archimedean,
  thm:isAtomicVectorLattice-iff-continuousPart-eq-zero}
Let $X$ be an Archimedean vector lattice. Then $X$ is atomic if and
only if the continuous part of $X$ equals $\{0\}$.
\end{theorem}

\section{Yudin's theorem}

The presence or absence of an infinite pairwise disjoint family of
non-zero vectors detects finite-dimensionality of an Archimedean
vector lattice. The main result of this section, due to Yudin,
characterises finite-dimensional Archimedean vector lattices as
those isomorphic, as vector lattices, to $\mathbb{R}^n$ for some
$n \in \mathbb{N}$.

\begin{theorem}
\label{thm:exists-disjoint-seq-of-not-finiteDimensional}
\lean{exists_disjoint_seq_of_not_finiteDimensional}
\leanok
\uses{def:isVLDisjoint, def:is-vl-archimedean, def:isVLAtom,
  def:isAtomicVectorLattice,
  thm:isAtomicVectorLattice-iff-forall-pos-dominates-atom,
  thm:isVLAtom-iff-no-disjoint-in-interval,
  thm:exists-projectionBand-principalBand-of-isVLAtom,
  thm:principalBand-eq-smul-of-isVLAtom}
Let $X$ be an Archimedean vector lattice that is not
finite-dimensional. Then there exists a pairwise disjoint sequence
$(x_n)_{n \in \mathbb{N}}$ of non-zero vectors in $X$.
\end{theorem}

\begin{theorem}
\label{thm:finiteDimensional-iff-latticeIso-pi}
\lean{finiteDimensional_iff_latticeIso_pi}
\leanok
\uses{def:is-vl-archimedean, def:isVLAtom, def:isAtomicVectorLattice,
  lem:pi-vectorLattice,
  thm:exists-disjoint-seq-of-not-finiteDimensional,
  thm:linearIndependent-of-pairwise-isVLDisjoint,
  thm:abs-sum-of-pairwise-isVLDisjoint,
  thm:isAtomicVectorLattice-iff-forall-pos-dominates-atom,
  thm:exists-projectionBand-principalBand-of-isVLAtom,
  thm:principalBand-eq-smul-of-isVLAtom}
\textbf{Yudin's theorem.} An Archimedean vector lattice $X$ is
finite-dimensional if and only if there exists $n \in \mathbb{N}$
such that $X$ is vector-lattice isomorphic to $\mathbb{R}^n$ (with
the pointwise lattice and vector-space structure).
\end{theorem}

\section{Coordinate expansion in atomic vector lattices}

In an atomic Archimedean vector lattice, every non-negative vector
admits a coordinate expansion along any maximal pairwise disjoint
family of atoms: it is the least upper bound of suitable scalar
multiples of the family.

\begin{theorem}
\label{thm:exists-isLUB-smul-of-isAtomic-maximal-disjoint}
\lean{exists_isLUB_smul_of_isAtomic_maximal_disjoint}
\leanok
\uses{def:isAtomicVectorLattice, def:is-vl-archimedean,
  def:isVLAtom, def:isMaximalDisjoint, def:isDisjointSet,
  thm:isAtomicVectorLattice-iff-forall-pos-dominates-atom,
  thm:isMaximalDisjoint-iff-forall-eq-zero,
  thm:exists-projectionBand-principalBand-of-isVLAtom,
  thm:principalBand-eq-smul-of-isVLAtom,
  thm:isVLDisjoint-or-smul-of-isVLAtom}
Let $X$ be an atomic Archimedean vector lattice and let $A \subseteq
X$ be a maximal pairwise disjoint subset of $X$ all of whose
elements are atoms. For every $x \in X$ with $0 \le x$ there exist
scalars $(c_a)_{a \in A}$ such that $x$ is the least upper bound of
$\{c_a\,a : a \in A\}$ in $X$.
\end{theorem}
